\section*{第 \ref{ch:b2:ism-star-formation} 章——星际介质与恒星形成}

\begin{solution}{exo:b2:ism-star-formation:1}
$n = \num{2.7e-24}/\num{1.67e-24} = 1.6$ 个氢原子每立方厘米。在空气中，
$n = P/k_{\mathrm B}T = 101325/(\num{1.381e-23}\times288) = \qty{2.55e25}{m^{-3}}$,
即每立方厘米 \num{2.5e19} 个分子，约多
$1.6\times10^{19}$ 倍。
\end{solution}

\begin{solution}{exo:b2:ism-star-formation:2}
$M_J \propto n^{-1/2}$：$n$ 除以 100，$M_J$ 乘以 10，因此
$M_J = 54.6$ 个太阳质量，此时粒子数密度为 $10^2$ 个粒子每立方厘米。云团包含
$\num{24000}/54.6 \approx 440$ 个\omterm{def:b2:ism-star-formation:jeans}{金斯质量}：它整体上不可能稳定，
必然会碎裂。
\end{solution}

\begin{solution}{exo:b2:ism-star-formation:3}
$c_s = \sqrt{\num{1.381e-23}\times10/(2.33\times\num{1.661e-27})} =
\qty{189}{m.s^{-1}}$（在 \qty{10}{K} 时）；在 \qty{100}{K} 时其为 $\sqrt{10}$ 倍，即 \qty{597}{m.s^{-1}}，
穿过 \qty{0.1}{pc} 所需时间为 $\num{3.086e15}/189 = \qty{1.63e13}{s}$，
即 \qty{0.52}{Ma}。
\end{solution}

\begin{solution}{exo:b2:ism-star-formation:4}
$\rho = 2.33\times\num{1.661e-27}\times n$，其中 $n$ 以 \unit{m^{-3}} 计。对于 $10^5$
个粒子每立方厘米，$\rho = \qty{3.87e-16}{kg.m^{-3}}$，且
$t_{\mathrm{ff}} = \sqrt{3\pi/(32G\rho)} = \qty{3.38e12}{s} = \qty{0.107}{Ma}$；对于
$10^2$ 个，$\rho$ 小 \num{1000} 倍，$t_{\mathrm{ff}}$ 则为 $\sqrt{1000} = 31.6$
倍，即 \qty{3.38}{Ma}。\omterm{def:b2:ism-star-formation:molecular-cloud}{致密云核}的坍缩快 32 倍。
\end{solution}

\begin{solution}{exo:b2:ism-star-formation:5}
由 $k_{\mathrm B}T/(\mu m_u) = c_s^2$，$(5k_{\mathrm B}T/G\mu m_u)^{3/2} =
5^{3/2}c_s^3G^{-3/2}$，由此得到该公式。由于 $M_J \propto \rho^{-1/2}$，
要使 $M_J$ 减半，密度需提高到四倍。由于 $M_J \propto T^{3/2}$，
将温度从 \qty{10}{K} 升至 \qty{20}{K}，使其乘以 $2^{3/2} = 2.83$。
\end{solution}

\begin{solution}{exo:b2:ism-star-formation:6}
$t_{\mathrm{KH}} = GM^2/(RL) = \num{6.674e-11}\times(\num{1.989e30})^2/(2\times
\num{6.96e8}\times2\times\num{3.83e26}) = \qty{2.48e14}{s}$，即
\qty{7.8}{Ma}。它比大多数盘的寿命长（半数盘在
约 \qty{3}{Ma} 后消失）：盘消失时恒星仍在收缩。
\end{solution}

\begin{solution}{exo:b2:ism-star-formation:7}
在 \qty{10}{K} 时，$c_s^3/G = 189^3/\num{6.674e-11} = \qty{1.01e17}{kg.s^{-1}}$，即
$\num{1.01e17}\times\num{3.156e7}/\num{1.989e30} = \num{1.6e-6}$ 个太阳质量每
年，一个太阳质量需 \qty{0.62}{Ma}。在 \qty{20}{K} 时，速率为
$2^{3/2} = 2.83$ 倍，即 \num{4.5e-6} 个太阳质量每年，一个太阳
质量需 \qty{0.22}{Ma}。
\end{solution}

\begin{solution}{exo:b2:ism-star-formation:8}
$r = \sqrt{L/(16\pi\sigma)}/T^2$，其中 $\sqrt{\num{3.83e26}/(16\pi\times\num{5.67e-8})}
= \qty{1.159e16}{m.K^2}$：在 \qty{1300}{K} 时，$r = \qty{6.86e9}{m} = \qty{0.046}{au}$，
远在水星轨道之内。在 \qty{5.2}{au} 处，$T = 278/\sqrt{5.2} = \qty{122}{K}$。
\end{solution}

\begin{solution}{exo:b2:ism-star-formation:9}
$v = \sqrt{GM/r} = \sqrt{\num{6.674e-11}\times1.3\times\num{1.989e30}/\num{1.496e13}}
= \qty{3.40}{km.s^{-1}}$，且 $P = 2\pi r/v = \qty{2.77e10}{s} = 877$ 年
（等价地为 $\sqrt{100^3/1.3}$ 年）。在一百万年内，气体已完成
$10^6/877 \approx \num{1140}$ 次旋转：与其自身的转动周期相比，这个盘很古老，
这也是其环保持光滑圆形的原因。
\end{solution}

\begin{solution}{exo:b2:ism-star-formation:10}
$M = \int_0^R 4\pi r^2\rho_0R/r\,\mathrm dr = 2\pi\rho_0R^3$，且在半径 $r$ 以内，
$m(r) = 2\pi\rho_0Rr^2$。由 $\mathrm dm = 4\pi\rho_0Rr\,\mathrm dr$，
\[
  U = -\int_0^R \frac{G\,m(r)}{r}\,\mathrm dm
    = -8\pi^2G\rho_0^2R^2\int_0^R r^2\,\mathrm dr = -\frac{8\pi^2}{3}G\rho_0^2R^5
    = -\frac23\,\frac{GM^2}{R},
\]
利用 $\rho_0 = M/(2\pi R^3)$。它是均匀情形值的 $\tfrac{2/3}{3/5} = 1.11$ 倍：
更为集中的云团束缚更紧，需要多 \qty{11}{\percent} 的热能
才能维持平衡；它更容易坍缩。
\end{solution}

\begin{solution}{exo:b2:ism-star-formation:11}
周期 $2\pi/\Omega = \qty{6.28e14}{s}$，约 \qty{20}{Ma}。取
$R = \qty{3.086e15}{m}$，$j = \Omega R^2 = \qty{9.52e16}{m^2.s^{-1}}$，且
$r_c = j^2/(GM) = \num{9.07e33}/\num{1.327e20} = \qty{6.83e13}{m} = \qty{457}{au}$。
在刚体转动中，$j = \Omega r^2$ 随离轴距离的平方增长：
来自核心内侧的气体角动量较小，会落在更靠近恒星的位置。
\end{solution}

\begin{solution}{exo:b2:ism-star-formation:12}
$\bar t = 2.46$, $\bar f = 0.588$; $\sum(t - \bar t)^2 = 12.612$,
$\sum(t - \bar t)(f - \bar f) = -2.018$，因此 $b = 2.018/12.612 = 0.160$ 每百万
年，且 $a = 0.588 + 0.160\times2.46 = 0.982$。该直线达到 $0.5$ 的时刻为
$(0.982 - 0.5)/0.160 = \qty{3.0}{Ma}$，降至零的时刻为 $0.982/0.160 = \qty{6.1}{Ma}$，如
图中所示。比例不可能低于零：较老的星团，例如 NGC 1960 在
\qty{30}{Ma} 时，仍会保持接近零的小比例，这是直线无法描述的。
\end{solution}

\begin{solution}{pb:b2:ism-star-formation:1}
\textbf{1.} $R = 0.08\times\num{3.086e16} = \qty{2.469e15}{m}$.

\textbf{2.} $N = \qty{1.5e26}{m^{-2}}$，因此 $2.8\times\num{1.661e-27}\times\num{1.5e26} =
\qty{0.698}{kg.m^{-2}}$.

\textbf{3.} $M = 0.698\times\pi R^2 = 0.698\times\num{1.915e31} = \qty{1.34e31}{kg}$,
即 $6.7$ 个太阳质量。

\textbf{4.} $V = \tfrac43\pi R^3 = \qty{6.30e46}{m^3}$，且
$\rho = M/V = \qty{2.12e-16}{kg.m^{-3}}$.

\textbf{5.} $n = \rho/(2.33m_u) = \qty{5.48e10}{m^{-3}}$，即 \num{5.5e4} 个粒子
每立方厘米。

\textbf{6.} 是的：其密度介于 $10^4$ 与 $10^5$ 之间，半径介于 0.03 与
\qty{0.2}{pc} 之间；其质量为 6.7 个太阳质量，处于云核质量范围的上限。

\textbf{7.} $M_J = 5.46\times(10^4/\num{5.48e4})^{1/2} = 2.33$ 个太阳质量。

\textbf{8.} $6.7/2.33 = 2.9$ 个\omterm{def:b2:ism-star-formation:jeans}{金斯质量}：云核不稳定，将发生垮塌。

\textbf{9.} $R_J = 0.1316\times(10^4/\num{5.48e4})^{1/2} = \qty{0.056}{pc}$，小于
半径 \qty{0.08}{pc}：结论相同。

\textbf{10.} $2^{3/2}\times2.33 = 6.6$ 个太阳质量：约等于云核质量。在
\qty{20}{K} 时云核将处于临界状态；正是低温使其发生垮塌。

\textbf{11.} $t_{\mathrm{ff}} = \sqrt{3\pi/(32\times\num{6.674e-11}\times\num{2.12e-16})}
= \qty{4.56e12}{s}$，约 \num{145000} 年。

\textbf{12.} $R/c_s = \num{2.469e15}/189 = \qty{1.31e13}{s}$，\num{414000} 年，是自由下落时间的三
倍。压力波无法在云核坍缩之前穿过云核；
压强的调整速度不够快，无法阻止坍缩。

\textbf{13.} $c_s = \qty{189}{m.s^{-1}}$（见\cref{exo:b2:ism-star-formation:3}）。

\textbf{14.} $c_s^3/G = \qty{1.01e17}{kg.s^{-1}}$，即每年 \num{1.6e-6} 个太阳质量。

\textbf{15.} $1/\num{1.6e-6} = \num{624000}$ 年。

\textbf{16.} 6.7 的三分之一约为 $2.2$：即两颗质量各为一个太阳质量的恒星。

\textbf{17.} $t_{\mathrm{KH}} = GM_\odot^2/(3R_\odot\times4L_\odot) = \qty{8.25e13}{s}$，
\qty{2.6}{Ma}。

\textbf{18.} 更长，为吸积时间的四倍：吸积停止时，恒星仍在收缩，
依靠收缩放光而未发生聚变；它是\omterm{def:b2:ism-star-formation:t-tauri}{金牛 T 型星}，尚未进入主
序。

\textbf{19.} $\Omega = \qty{0.4e3}{m.s^{-1}}/\qty{3.086e16}{m} = \qty{1.30e-14}{rad.s^{-1}}$.

\textbf{20.} $j = 1.30\times10^{-14}\times(\num{2.469e15})^2 = \qty{7.90e16}{m^2.s^{-1}}$.

\textbf{21.} $r_c = j^2/(GM_\odot) = \num{6.24e33}/\num{1.327e20} = \qty{4.70e13}{m}$，即
\qty{314}{au}：大于 HL 金牛盘，
符合来自云核边缘气体的预期。

\textbf{22.} $T = (L/16\pi\sigma r^2)^{1/4} = \qty{278}{K}$，位于 \qty{1}{au} 处，而
$278/\sqrt{30} = \qty{51}{K}$，位于 \qty{30}{au} 处。

\textbf{23.} $r = \num{1.159e16}/180^2 = \qty{3.58e11}{m}$，即 \qty{2.4}{au}，与正文一致。

\textbf{24.} 一个太阳质量的恒星需要
\textbf{\qty[text-series-to-math=true]{624}{} 千年}才能形成，约为
其云核自由下落时间（\num{145000} 年）的四倍：外层气体只能逐步到达
\omterm{def:b2:ism-star-formation:protostar}{原恒星}，且自内向外进行。
\end{solution}
